% chktex-file 8
% chktex-file 12
% chktex-file 18
\beginsong{Točí se točí}[by={Petr Lutka}]

\beginverse{}
\ch{C}{}{}{}Točí se točí \ch{G}{}{}{}kolo ve mlý\ch{C}{}{}{}ně.
\ch{C}{}{}{}Točí se točí já \ch{F}{}{}{}jsem v tom nevi\ch{C}{}{}{}nně.
Točí se \ch{G}{}{}{}kolo voda vzduchem \ch{C}{}{}{}padá
\ch{C}{}{}{}já tomu mlejnu \ch{G}{}{}{}ukazuju \ch{C}{}{}{}záda.
\endverse{}

\beginverse{}
V tom mlejně žila překrásná dívka
kterou jsem měl tuze moc rád,
měla modrý oči a černý vlasy
jenom jednu chybu: nechtěla mi dát.
\endverse{}

\beginverse{}
Nechtěla mi dát ani políbení
nechtěla mi dát ani hubičku.
Proto jsem jí musel, proto jsem jí musel,
proto jsem jí musel utopiti v rybníčku.
\endverse{}

\beginverse{}
Když se dcera domů dlouho nevracela
její otec mlynář o ní starost měl,
poslal za ní mládka, kouknout se co dělá
proto poslal mládka, aby za ní šel.
\endverse{}

\beginverse{}
Když tam mládek přišel, právě topila se,
její ručka bílá o pomoc prosila,
zachránit ji nemoh, kosa byla ostrá
jeho mladý čelo správně trefila.
\endverse{}

\beginverse{}
Když se mládek s dcerou dlouho nevraceli,
přišel se sám mlynář, na dcerušku ptát.
Když kameny mlýnský jeho tělo mlely,
řekl jsem si: musíš na cestu se dát.
\endverse{}

\beginverse{}
= 1.
\endverse{}

\endsong{}
